\section{Calibração Espacial por Homografia Métrica em Vídeos não Ortogonais}

Em vídeos gravados com a câmera inclinada em relação ao plano do movimento, a distância medida diretamente em pixels não é uniforme em toda a imagem. Um mesmo comprimento físico pode ocupar quantidades diferentes de pixels dependendo da posição na cena. Por isso, a calibração por apenas dois pontos é adequada apenas quando a câmera está aproximadamente perpendicular ao plano analisado.

Quando o cenário contém linhas do piso, marcas de referência ou outro retângulo físico conhecido, é possível usar uma homografia para transformar a imagem oblíqua em uma vista retificada do plano. Considere quatro pontos da imagem:

\[
    \mathbf{p}_1,\mathbf{p}_2,\mathbf{p}_3,\mathbf{p}_4,
\]

marcados na ordem superior esquerdo, superior direito, inferior direito e inferior esquerdo. Esses pontos devem corresponder aos cantos de um retângulo real no plano do movimento, com largura $L$ e altura $A$.

A homografia é uma transformação projetiva representada por uma matriz $H \in \mathbb{R}^{3\times3}$, definida a menos de um fator escalar:

\[
    \lambda
    \begin{bmatrix}
    x'\\y'\\1
    \end{bmatrix}
    =
    H
    \begin{bmatrix}
    x\\y\\1
    \end{bmatrix}.
\]

No sistema implementado, os quatro pontos da imagem são associados aos quatro pontos de destino:

\[
    (0,0),\quad
    (L\rho,0),\quad
    (L\rho,A\rho),\quad
    (0,A\rho),
\]

onde $\rho$ é a resolução escolhida para a planta retificada, em pixels por unidade de medida. Assim, a imagem resultante passa a representar uma vista superior aproximada do plano calibrado, com escala espacial conhecida:

\[
    S = \frac{1}{\rho}.
\]

Na prática, por causa do arredondamento da dimensão final em pixels, a escala usada pelo aplicativo é registrada como:

\[
    S = \frac{L}{N_x-1},
\]

em que $N_x$ é a largura, em pixels, da imagem retificada. A origem pode ser fixada no canto inferior esquerdo do retângulo retificado. Desse modo, para um ponto rastreado com coordenadas de imagem $(p_x,p_y)$ na imagem corrigida, obtém-se:

\[
    x = p_x S,
    \qquad
    y = (N_y-p_y)S,
\]

onde $N_y$ é a altura da imagem retificada.

\subsection{Interpretação}

Esse procedimento separa a limitação experimental da limitação matemática. A homografia não ``corrige'' qualquer vídeo de modo absoluto; ela corrige apenas pontos pertencentes ao mesmo plano usado na calibração. Portanto, a técnica é adequada quando a trajetória analisada pode ser aproximada como pertencente ao plano definido pelas referências marcadas, ou quando o erro de profundidade é pequeno em relação aos objetivos didáticos da atividade.

\subsection{Limitações}

As principais limitações são:

\begin{itemize}
    \item os quatro pontos devem pertencer ao mesmo plano físico;
    \item a câmera deve permanecer fixa durante todo o vídeo;
    \item o objeto analisado deve se mover aproximadamente no plano calibrado;
    \item marcações imprecisas dos cantos propagam erro para a escala;
    \item objetos com deslocamento significativo em profundidade não são completamente corrigidos por uma única homografia.
\end{itemize}

Apesar dessas restrições, o uso de referências coplanares amplia a robustez do sistema em vídeos não ortogonais, pois substitui uma calibração local por dois pontos por uma calibração global do plano de movimento. Isso torna a escala espacial mais coerente em toda a região retificada e melhora a interpretação dos gráficos cinemáticos obtidos.
